\begin{table}[tb]
\centering\small
\caption{Resolving ties among the same highest-vote candidates with a forest for masking and refitting. Each entry gives effective compression and outer breach rate ($U$/breach, \%). Order selects the earliest eligible candidate in the frozen slate order. Uniform averages the guarded outcomes of every tied nominee. The local rows average the three draws within each fold.}
\label{tab:within-tier-forest}
\begin{tabular}{@{}llrrrr@{}}
\toprule
Slate & Panel & Order & Uniform & Point & Full\\
\midrule
Gain & Ten  & 2.89/3.33   & 5.42/5.97   & 6.33/6.67   & 6.33/10.00\\
     & Wide & 8.28/16.67  & 11.06/13.89 & 13.77/16.67 & 13.77/16.67\\
Local & Ten  & 18.65/0.00 & 16.68/0.97  & 17.51/1.11  & 16.43/2.22\\
      & Wide & 11.57/3.70 & 11.63/6.59  & 12.46/3.70  & 12.89/5.56\\
\bottomrule
\end{tabular}
\end{table}

\paragraph{How much do the margins add within a tie?}
On the original ten tasks, replacing the fixed order by uniform choice contributed 2.53 of the 3.44 percentage points gained by using the full margin. The remaining 0.91 points (paired 95\% interval $[-2.89,5.63]$) came from Dry Bean and Steel, with within-task differences of $+18.75$ and $-9.62$ points. The full margin also raised the breach rate by 4.03 points over uniform. On the wider Gain slates uniform choice contributed 2.78 of the 5.49 points. On the local slates full minus uniform was $-0.25$ and $+1.26$ points. All four intervals for full minus uniform $U$ included zero (Table~\ref{tab:within-tier-forest-contrasts}). These primary means include all folds, including those without a tie. On the wider Gain panel, restricting the comparison to visible ties gave $+9.50$ points for full minus uniform ($[-2.77,23.17]$), compared with $+2.71$ across all folds. Only 5 of 30 original ten-task folds and 5 of 16 wider folds had tied candidates with different realized guarded $U$.
